%------------------------------------------------------------------------------%
% \chapter{Revisão}
%------------------------------------------------------------------------------%

%------------------------------------------------------------------------------%
\begin{exercicios}{Conjuntos Numéricos e Desigualdades}
    
    \begin{exercicio}
        Determine se as afirmações são verdadeiras ou falsas.
        \begin{mathlist}[2]
            \item -3 < -20
            \item  \frac{2}{3} >  \frac{5}{6}
            \item -\frac{5}{6} < -\frac{11}{12}
        \end{mathlist}
        \begin{answers}
            \begin{alphalist}[3]
            \item Falso
            \item Falso
            \item Falso
        \end{alphalist}
        \end{answers}
    \end{exercicio}
        
    \begin{exercicio}
        Represente os intervalos dados na reta real.
        \begin{mathlist}[2]
            \item (3, 6)
            \item [-1, 4)
            \item \left[ - \frac{6}{5}, - \frac{1}{2} \right] 
            \item (- \infty, 5]
        \end{mathlist}
    \end{exercicio}
        
    \begin{exercicio}
        Encontre os valores de $x$ que satisfazem as desigualdades.
        \begin{mathlist}[2]
            \item 2x + 4 < 8
            \item -6 > 4 + 5x
            \item -4x \geq 20
            \item -6 < x - 2 < 4
            \item 0 \leq x + 1 \leq 4
            \item x + 1 > 4 \) ou $x + 2 < -1$
            \item x - 4 \leq 1 \) e $x + 3 > 2$
        \end{mathlist}
    \end{exercicio}
    
    \begin{exercicio}
        Descreva os intervalos abaixo usando desigualdades e desenhe-os na reta real.
        \begin{mathlist}[2]
            \item (-2,0)
            \item [1,6)
            \item (-3,\infty)
            \item (-\infty; 12,5]
            \item [-4,5]
            \item (-5, -2]
        \end{mathlist}    
    \end{exercicio}
    
    \begin{exercicio}
        Escreva os conjuntos abaixo na forma de intervalos e desenhe-os na reta real.
        \begin{mathlist}
            \item \{ x \in \R \st x \geq 0,17 \}
            \item \{ x \in \R \st x \geq 4 \}
            \item \{ x \in \R \st - 3 < x < -1 \}
            \item \{ x \in \R \st - 1 \leq x \leq 0 \}
            \item \{ x \in \R \st 1/100 \leq x < 100 \} 
            \item \{ x \in \R \st x \leq -2 \text{ ou } x > 5 \}
            \item \{ x \in \R \st x > -4 \}
            \item \{ x \in \R \st x < 3 \}
        \end{mathlist}
    \end{exercicio}
    
    \begin{exercicio}
        Considerando os conjuntos
        \[ A = \{x \in \R \st x \geq 1 \} \]
        \[ B = \{x \in \R \st x < 2 \} \]
        \[ C = \{x \in \R \st - 2 < x \leq 4 \} \]
        determine
        \begin{mathlist}[2]
            \item A\cup C
            \item B\cup C
            \item A\cup B
            \item A\cap C
            \item B\cap C
            \item A\cap B
        \end{mathlist}
    \end{exercicio}

    \begin{exercicio}
        Reescreva os intervalos abaixo na forma mais simples e compacta possível.
        \begin{mathlist}[2]
            \item [-4,5] \cup [-2,1]
            \item [-2,3] \cup [3,4]
            \item (0,2) \cup [2,8]
            \item (-1,4) \cup (4,6)
            \item [-4,1) \cup (-2,2]
            \item [-4, -52) \cup (-3, -2)
            \item (-\infty,6) \cup (4,7)
            \item (-\infty,5) \cup (-1,\infty)
        \end{mathlist}
    \end{exercicio}
    
    \begin{exercicio}
        Escreva os conjuntos abaixo usando desigualdades e represente-os na reta real.
        \begin{mathlist}[2]
            \item (-3,1) \cup (-1,2)
            \item [-2,2) \cap ( 1/2,4]
            \item [ 1,4) \cup (1,6]
            \item (-\infty, 2] \cap (-2,0]
            \item (-\infty,-2] \cup [3,\infty)
            \item (-\infty, 8) \cap (8,\infty)
        \end{mathlist}
    \end{exercicio}
    
    \begin{exercicio}
        Descreva os conjuntos abaixo usando a notação de intervalo.
        \addgraphic{.8}{exercicio-intervalos-01}
    \end{exercicio}
    
    \begin{exercicio}
        Considerando os conjuntos
        \[
        A = (-\infty, -3] \hfill
        B = (-1,7)        \hfill
        C = [-5, 6]       \hfill\null 
        \]
        determine
        \begin{mathlist}[2]
            \item  A \cup C
            \item  B \cup C
            \item  A \cap C
            \item  B \cap C
            \item  A \cup B \cup C
            \item  A \cap B \cap C
            \item  (A \cup B) \cap C
            \item  A \cup (B \cap C)
            \item  (A \cap B) \cup C
        \end{mathlist}
    \end{exercicio}
        
    

% 1
\begin{exercicio}
    Resolva as inequações.
    \begin{mathlist}[2]
        \item 2x > 3
        \item 8x \geq -5
        \item x - 4 \leq 5
        \item \frac{a}{2} < 7
        \item 3z - \frac{1}{2} > \frac{1}{4}
        \item x + 1 \geq -1
        \item -x \leq 6
        \item 3 \geq -9x
        \item -\frac{w}{4} > 8
        \item 3 -2y < 7
    \end{mathlist}
    \begin{answers}
        \begin{mathlist}[3]
            \item x \in \left(\sfrac{3}{2}, \infty\right)
            \item x \in \left[- \sfrac{5}{8}, \infty\right)
            \item x \in \left(-\infty, 9\right]
            \item a \in \left(-\infty, 14\right)
            \item z \in \left(\sfrac{1}{4}, \infty\right)
            \item x \in \left[-2, \infty\right)
            \item x \in \left[-6, \infty\right)
            \item x \in \left[- \sfrac{1}{3}, \infty\right)
            \item w \in \left(-\infty, -32\right)
            \item y \in \left(-2, \infty\right)
        \end{mathlist}
    \end{answers}
\end{exercicio}

% 2
\begin{exercicio}
  Resolva as inequações.
  \begin{mathlist}[1]
    \item 1 - 2(x - 1) < 2
    \item 2 - 3x \geq x + 14
    \item 5v - 32 \leq 4 - 7v
    \item 2 - z > 3(z + 3)
    \item 2(3x + 1) < 4(5 - 2x)
    \item 8(x + 3) > 12(1 - x)
    \item \frac{2}{3} - \frac{1}{2} x \geq \frac{1}{6} + x
    \item 3(3x - 2) + 2 \left(x + \frac{1}{2} \right) \leq 19 - x
    \item \frac{3x}{2} + \frac{x}{3} + \frac{x}{6} > 0
    \item \frac{1}{3} + \frac{x}{2} < \frac{5}{6} - \frac{2x}{3}
    \item \frac{3x+1}{4} - 1 \geq \frac{1}{2} - 2x
    \item \frac{1-2x}{3} + \frac{x-2}{6} > \frac{x+3}{2} - 1
    \item \frac{2}{5}x + 1 \leq \frac{1}{5} - 2x
    \item \frac{x+2}{3} + \frac{2-3x}{2} < \frac{4x}{3}
    \item \frac{x}{3} - \frac{x+1}{2} < \frac{1-x}{4}
    \item 3(1 - 2x) < 2(x + 1) + x - 7
    \item \frac{x+10}{5} > -x + 6
    \item \frac{3x-1}{4} + \frac{1-4x}{2} < 1
    \item \frac{2x}{5} + 1 \leq 2 \left(x + \frac{3}{5} \right)
    \item \frac{1-2x}{3} - \frac{1+3x}{2} \geq 2
  \end{mathlist}
  \begin{answers}
      \begin{mathlist}[3]
        \item x \in \left(\sfrac{1}{2}, \infty\right)
        \item x \in \left(-\infty, -3\right]
        \item v \in \left(-\infty, 3\right]
        \item z \in \left(-\infty, - \sfrac{7}{4}\right)
        \item x \in \left(-\infty, \sfrac{9}{7}\right)
        \item x \in \left(- \sfrac{3}{5}, \infty\right)
        \item x \in \left(-\infty, \sfrac{1}{3}\right]
        \item x \in \left(-\infty, 2\right]
        \item x \in \left(0, \infty\right)
        \item x \in \left(-\infty, \sfrac{3}{7}\right)
        \item x \in \left[\sfrac{5}{11}, \infty\right)
        \item x \in \left(-\infty, - \sfrac{1}{2}\right)
        \item x \in \left(-\infty, - \sfrac{1}{3}\right]
        \item x \in \left(\sfrac{2}{3}, \infty\right)
        \item x \in \left(-\infty, 9\right)
        \item x \in \left(\sfrac{8}{9}, \infty\right)
        \item x \in \left(\sfrac{10}{3}, \infty\right)
        \item x \in \left(- \sfrac{3}{5}, \infty\right)
        \item x \in \left[- \sfrac{1}{8}, \infty\right)
        \item x \in \left(-\infty, -1\right]
    \end{mathlist}
  \end{answers}
\end{exercicio}

% 3
\begin{exercicio}
    Resolva as inequações e represente suas soluções na reta real.
    \begin{mathlist}[2]
        \item 1 < 2x < 3
        \item -3 \leq 4x \leq 8
        \item -1 \leq x + 2 \leq 5
        \item 0 \leq 2x - 2 \leq 6
        \item -6 \leq -2(x - 1) \leq 0
        \item 2 \leq \frac{x}{3} < 4
        \item -3 < \frac{3x}{2} \leq 6
        \item -\frac{1}{4} \leq \frac{3x-4}{7} \leq \frac{1}{2}
        \item \frac{1}{6} < \frac{2x-13}{12} < \frac{2}{3}
        \item -6 \leq 15-3(4x+7) \leq 18
        \item -4 \leq \frac{5x-4}{6} \leq 1
        \item -1 \leq \frac{4s-3}{5} \leq \frac{3}{15}
        \item -9 \leq \frac{5-8x}{3} \leq 1
        \item -\frac{5}{4} \leq \frac{2x-3}{2} \leq \frac{7}{2}
        \item -\frac{3}{2} \leq \frac{2x-5}{3} \leq \frac{1}{6}
    \end{mathlist}
    \begin{answers}
        \begin{mathlist}[3]
            \item x \in \left(\sfrac{1}{2}, \sfrac{3}{2}\right)
            \item x \in \left[- \sfrac{3}{4}, 2\right]
            \item x \in \left[-3, 3\right]
            \item x \in \left[1, 4\right]
            \item x \in \left[1, 4\right]
            \item x \in \left[6, 12\right)
            \item x \in \left(-2, 4\right]
            \item x \in \left[\sfrac{3}{4}, \sfrac{5}{2}\right]
            \item x \in \left(\sfrac{15}{2}, \sfrac{21}{2}\right)
            \item x \in \left[-2, 0\right]
            \item x \in \left[-4, 2\right]
            \item s \in \left[- \sfrac{1}{2}, 1\right]
            \item x \in \left[\sfrac{1}{4}, 4\right]
            \item x \in \left[\sfrac{1}{4}, 5\right]
            \item x \in \left[\sfrac{1}{4}, \sfrac{11}{4}\right]
        \end{mathlist}
    \end{answers}
\end{exercicio}

% 4
\begin{exercicio}
    Resolva as desigualdades e represente a solução na reta real.
    \begin{mathlist}[2]
        \item \left(x - 2\right) \left(x - 4\right) \geq 0
        \item \left(x + 1\right) \left(x - 3\right) \leq 0
        \item x \left(2 x - 1\right) \geq 0
        \item 2 x \left(x - \frac{1}{4}\right) \leq 0
        \item -3 \left(x + 2\right) \left(x - 3\right) < 0
        \item \left(x + 3\right) \left(3-5x\right) \geq 0
        \item \left(x - \frac{1}{2}\right) \left(2 x + 5\right) \leq 0
        \item \left(x - 6\right)^{2} > 0
    \end{mathlist}
    \begin{answers}
        \begin{mathlist}[2]
            \item x \in \left(-\infty, 2\right] \cup \left[4, \infty\right)
            \item x \in \left[-1, 3\right]
            \item x \in \left(-\infty, 0\right] \cup \left[\sfrac{1}{2}, \infty\right)
            \item x \in \left[0, \sfrac{1}{4}\right]
            \item x \in \left(-\infty, -2\right) \cup \left(3, \infty\right)
            \item x \in \left[-3, \sfrac{3}{5}\right]
            \item x \in \left[- \sfrac{5}{2}, \sfrac{1}{2}\right]
            \item x \in \left(-\infty, 6\right) \cup \left(6, \infty\right)
        \end{mathlist}
    \end{answers}
\end{exercicio}

% 5
\begin{exercicio}
    Resolva as desigualdades e represente a solução na reta real.
    \begin{mathlist}[2]
        \item x^2 - 3x \geq 0
        \item 3x^2 \leq 5x
        \item x^2 - 8 \leq 0
        \item x^2 + 6x \leq 0
        \item 3x^2 - \sqrt{5\,}x \geq 0
        \item x^2 + 2x > 3
        \item 49x^2 \leq 9
        \item -x^2 + 5 \leq 0
        \item -2x^2 + x \geq -6
        \item x^2 + 4x + 7 \leq 0
        \item x^2 + 2x + 1 \leq 0
        \item 2x^2 \geq 20 - 6x
        \item x^2 + 9x + 18 \leq 0
        \item x^2 - 6x + 9 \geq 0
        \item -3x^2 + 16x \leq 5
        \item 16x^2 + 25 \leq 0
        \item -4x^2 + 12x - 9 \leq 0
        \item 3x^2 \leq 2x + 5
        \item -2x^2 + 8x + 24 \leq 0
        \item -x^2 + 20x - 36 \geq 0
        \item 2x^2 - 5x \geq 3
        \item (x - 6)^2 \geq 4
    \end{mathlist}
    \begin{answers}
        \begin{mathlist}[2]
            \item x \in \left(-\infty, 0\right] \cup \left[3, \infty\right)
            \item x \in \left[0, \sfrac{5}{3}\right]
            \item x \in \left[- 2 \sqrt{2}, 2 \sqrt{2}\right]
            \item x \in \left[-6, 0\right]
            \item x \in \left(-\infty, 0\right] \cup \left[\sfrac{\sqrt{5}}{3}, \infty\right)
            \item x \in \left(-\infty, -3\right) \cup \left(1, \infty\right)
            \item x \in \left[- \sfrac{3}{7}, \sfrac{3}{7}\right]
            \item x \in \left(-\infty, - \sqrt{5}\right] \cup \left[\sqrt{5}, \infty\right)
            \item x \in \left[- \sfrac{3}{2}, 2\right]
            \item x \in \emptyset
            \item x \in \left\{-1\right\}
            \item x \in \left(-\infty, -5\right] \cup \left[2, \infty\right)
            \item x \in \left[-6, -3\right]
            \item x \in \left(-\infty, \infty\right)
            \item x \in \left(-\infty, \sfrac{1}{3}\right] \cup \left[5, \infty\right)
            \item x \in \emptyset
            \item x \in \left(-\infty, \infty\right)
            \item x \in \left[-1, \sfrac{5}{3}\right]
            \item x \in \left(-\infty, -2\right] \cup \left[6, \infty\right)
            \item x \in \left[2, 18\right]
            \item x \in \left(-\infty, - \sfrac{1}{2}\right] \cup \left[3, \infty\right)
            \item x \in \left(-\infty, 4\right] \cup \left[8, \infty\right)
        \end{mathlist}
    \end{answers}
\end{exercicio}

% 6 
\begin{exercicio}
    Resolva as desigualdades e represente a solução na reta real.
    \begin{mathlist}[2]
        \item  1 \leq x^2 + 2x - 2 \leq 6
        \item -4 \leq 3x^2 - 10 \leq 2
        \item -3 \leq x^2 - 4x \leq 5
        \item -2 \leq 2x^2 + 3x + 4 \leq 3
        \item  4 \leq (x - 6)^2 \leq 9
        \item  0 \leq x^2 - x \leq 20
    \end{mathlist}
    \begin{answers}
        \begin{mathlist}[2]
            \item x \in \left[-4, -3\right] \cup \left[1, 2\right]
            \item x \in \left[-2, - \sqrt{2}\right] \cup \left[\sqrt{2}, 2\right]
            \item x \in \left[-1, 1\right] \cup \left[3, 5\right]
            \item x \in \left[-1, - \sfrac{1}{2}\right]
            \item x \in \left[3, 4\right] \cup \left[8, 9\right]
            \item x \in \left[-4, 0\right] \cup \left[1, 5\right]
        \end{mathlist}
    \end{answers}
\end{exercicio}

% 7 
\begin{exercicio}
    Resolva as inequações e represente a solução na reta real.
    \begin{mathlist}[2]
        \item \frac{x-2}{x+3}    \leq 0
        \item \frac{x+4}{x-2}    \geq 0
        \item \frac{2x-3}{x-1}   \leq 0
        \item \frac{4x+5}{x+2}   \geq 0
        \item \frac{x-3}{2x+6}   \leq 0
        \item \frac{4-5x}{2x-1}  \geq 0
        \item 3 - \frac{x}{x+2}  \leq 0
        \item \frac{3x-2}{5-2x}  \geq 0
        \item \frac{5x}{x-4}     \geq 10
        \item \frac{2x-7}{x-2}   \geq 3
        \item \frac{3x+10}{2x-5} \geq -3
        \item \frac{3x-4}{1-6x}  \leq 2
        \item \frac{6-x}{x-4}    \geq 1
    \end{mathlist}
    \begin{answers}
        \begin{mathlist}[2]
            \item x \in \left(-3, 2\right]
            \item x \in \left(-\infty, -4\right] \cup \left(2, \infty\right)
            \item x \in \left(1, \sfrac{3}{2}\right]
            \item x \in \left(-\infty, -2\right) \cup \left[- \sfrac{5}{4}, \infty\right)
            \item x \in \left(-3, 3\right]
            \item x \in \left(\sfrac{1}{2}, \sfrac{4}{5}\right]
            \item x \in \left[-3, -2\right)
            \item x \in \left[\sfrac{2}{3}, \sfrac{5}{2}\right)
            \item x \in \left(4, 8\right]
            \item x \in \left[-1, 2\right)
            \item x \in \left(-\infty, \sfrac{5}{9}\right] \cup \left(\sfrac{5}{2}, \infty\right)
            \item x \in \left(-\infty, \sfrac{1}{6}\right) \cup \left[\sfrac{2}{5}, \infty\right)
            \item x \in \left(4, 5\right]
        \end{mathlist}
    \end{answers}
\end{exercicio}

% 8 
\begin{exercicio}
    Resolva as inequações e represente a solução na reta real.
    \begin{mathlist}[2]
        \item \frac{x^2 -3x+2}{x^2 +x} \geq 0
        \item \frac{x+8}{x^2 +7x+12} \leq 0
        \item \frac{x^2 +5}{x^2 -4} \geq 5
        \item \frac{x^2 +6}{x^2 +1} \leq 2
        \item 1 + \frac{2}{x+1} \leq \dfrac{2}{x}
        \item \frac{4-2x}{x^2 -4} \leq 3
        \item \frac{x+6}{3x^2 +2} \geq 1
        \item \frac{x-5}{2x-5} \geq x
        \item \frac{4x-7}{x+2} \leq x - 2
        \item \frac{3x-1}{x+4} + \dfrac{x}{x-5} \leq 0
        \item \frac{x}{x+1} - \dfrac{1}{x-3} \geq 2
        \item \frac{x^2 +2x+3}{x+15} \leq 1
        \item \frac{3x^2 +2x-13}{x^2 -3x-10} \geq 2
    \end{mathlist}
    \begin{answers}
        \begin{mathlist}[2]
            \item x \in \left(-\infty, -1\right) \cup \left(0, 1\right] \cup \left[2, \infty\right)
            \item x \in \left(-\infty, -8\right] \cup \left(-4, -3\right)
            \item x \in \left[- \sfrac{5}{2}, -2\right) \cup \left(2, \sfrac{5}{2}\right]
            \item x \in \left(-\infty, -2\right] \cup \left[2, \infty\right)
            \item x \in \left[-2, -1\right) \cup \left(0, 1\right]
            \item x \in \left(-\infty, - \sfrac{8}{3}\right] \cup \left(-2, 2\right) \cup \left(2, \infty\right)
            \item x \in \left[-1, \sfrac{4}{3}\right]
            \item x \in \left(-\infty, \sfrac{5}{2}\right)
            \item x \in \left(-2, 1\right] \cup \left[3, \infty\right)
            \item x \in \left(-4, \sfrac{1}{2}\right] \cup \left[\sfrac{5}{2}, 5\right)
            \item x \in \left[- \sqrt{5}, -1\right) \cup \left[\sqrt{5}, 3\right)
            \item x \in \left(-\infty, -15\right) \cup \left[-4, 3\right]
            \item x \in \left(-\infty, -7\right] \cup \left(-2, -1\right] \cup \left(5, \infty\right)
        \end{mathlist}
    \end{answers}
\end{exercicio}

% 9 
\begin{exercicio}
    Resolva as inequações e represente a solução na reta real.
    \begin{mathlist}[2]
        \item \sqrt{x} - 8              \leq 0
        \item \sqrt{x} - 3              \geq 0
        \item \sqrt{x} + 10             \geq 0
        \item \sqrt{6 - 5x} - 4         \geq 0
        \item \sqrt{2x - 3}             \leq 5
        \item \sqrt{3x + 12} + 7        \geq 0
        \item \sqrt{x + 2}              \geq x - 4
        \item \sqrt{10 - 3x}            \leq x - 2
        \item \sqrt{8x + 9}             \leq 2x + 1
        \item \sqrt{9x^2 - 1}           \geq 2 - 3x
        \item \sqrt{8 - 4x^2} + x       \leq 3
        \item \sqrt{x^2 - 4x - 12}      \leq 3x + 2
        \item \sqrt{2x^2 + 1} - 2x + 1  \geq 0
        \item \sqrt{-x^2 + 5x + 14} - 2 \leq x
    \end{mathlist}
    \begin{answers}
        \begin{mathlist}[2]
            \item x \in \left[0, 64\right]
            \item x \in \left[9, \infty\right)
            \item x \in \left[0, \infty\right)
            \item x \in \left(-\infty, -2\right]
            \item x \in \left[\sfrac{3}{2}, 14\right]
            \item x \in \left[-4, \infty\right)
            \item x \in \left[-2, 7\right]
            \item x \in \left[3, \sfrac{10}{3}\right]
            \item x \in \left[2, \infty\right)
            \item x \in \left[\sfrac{5}{12}, \infty\right)
            \item x \in \left[- \sqrt{2}, \sfrac{1}{5}\right] \cup \left[1, \sqrt{2}\right]
            \item x \in \left[6, \infty\right)
            \item x \in \left(-\infty, 2\right]
            \item x \in \left\{-2\right\} \cup \left[\sfrac{5}{2}, 7\right]
        \end{mathlist}
    \end{answers}
\end{exercicio}


%    \begin{exercicio}
%        Resolva as inequações e represente suas soluções na reta real
%        \begin{mathlist}[2]
%            \item 1 < 2x < 3
%            \item -3 \leq 4x \leq 8
%            \item -1 \leq x + 2 \leq 5
%            \item 0 \leq 2x - 2 \leq 6
%            \item -6 \leq -2(x - 1) \leq 0
%            \item 2 \leq \frac{x}{3} < 4
%            \item -3 < \frac{3x}{2} \leq 6
%            \item -\frac{1}{4} \leq \frac{3x-4}{7} \leq \frac{1}{2}
%            \item \frac{1}{6} < \frac{2x-13}{12} < \frac{2}{3}
%            \item -6 \leq 15-3(4x+7) \leq 18
%            \item -4 \leq \frac{5x-4}{6} \leq 1
%            \item -1 \leq \frac{4s-3}{5} \leq \frac{3}{15}
%            \item -9 \leq \frac{5-8x}{3} \leq 1
%            \item -\frac{5}{4} \leq \frac{2x-3}{2} \leq \frac{7}{2}
%            \item -\frac{3}{2} \leq \frac{2x-5}{3} \leq \frac{1}{6}
%        \end{mathlist}
%        \todo{Copiar resposta do exercício 4 da pg 70 do arquivo 2 do livro do Francisco}
%    \end{exercicio}                
%                               
%    \begin{exercicio}
%        Resolva as desigualdades e represente a solução na reta real
%        \begin{mathlist}[2]
%            \item (x - 2)(x - 4) \geq 0
%            \item (x + 1)(x - 3) \leq 0
%            \item (2x - 1)x \geq 0
%            \item 2x(x - 1/4) \leq 0
%            \item -3(x + 2)(x - 3) < 0
%            \item (3 - 5x)(x + 3) \geq 0
%            \item (2x + 5) (x - 1/2 ) \leq 0
%            \item (x - 6)^2 > 0
%        \end{mathlist}
%        \todo{Copiar resposta do exercício 2 da pg 105 do arquivo 2 do livro do Francisco}
%    \end{exercicio}
%
%    \begin{exercicio}
%        Resolva as desigualdades e represente a solução na reta real
%        \begin{mathlist}[2]
%            \item x^2 - 3x \geq 0
%            \item 3x^2 \leq 5x
%            \item x^2 - 8 \leq 0
%            \item x^2 + 6x \leq 0
%            \item 3x^2 - \sqrt{5\,}x \geq 0
%            \item x^2 + 2x > 3
%            \item 49x^2 \leq 9
%            \item -x^2 + 5 \leq 0
%            \item -2x^2 + x \geq -6
%            \item x^2 + 4x + 7 \leq 0
%            \item x^2 + 2x + 1 \leq 0
%            \item 2x^2 \geq 20 - 6x
%            \item x^2 + 9x + 18 \leq 0
%            \item x^2 - 6x + 9 \geq 0
%            \item -3x^2 + 16x \leq 5
%            \item 16x^2 + 25 \leq 0
%            \item -4x^2 + 12x - 9 \leq 0
%            \item 3x^2 \leq 2x + 5
%            \item -2x^2 + 8x + 24 \leq 0
%            \item -x^2 + 20x - 36 \geq 0
%            \item 2x^2 - 5x \geq 3
%            \item (x - 6)^2 \geq 4
%        \end{mathlist}
%        \todo{Copiar resposta do exercício 4 da pg 105 do arquivo 2 do livro do Francisco}
%    \end{exercicio}
%
%    \begin{exercicio}
%        Resolva as desigualdades e represente a solução na reta real
%        \begin{mathlist}[2]
%            \item  1 \leq x^2 + 2x - 2 \leq 6
%            \item -4 \leq 3x^2 - 10 \leq 2
%            \item -3 \leq x^2 - 4x \leq 5
%            \item -2 \leq 2x^2 + 3x + 4 \leq 3
%            \item  4 \leq (x - 6)^2 \leq 9
%            \item  0 \leq x^2 - x \leq 20
%        \end{mathlist}
%        \todo{Copiar resposta do exercício 5 da pg 105 do arquivo 2 do livro do Francisco}
%    \end{exercicio}
%
%    \begin{exercicio}
%        Resolva as inequações e represente a solução na reta real
%        \begin{spacing}{2}
%            \begin{mathlist}[2]
%                \item \dfrac{x-2}{x+3}    \leq 0
%                \item \dfrac{x+4}{x-2}    \geq 0
%                \item \dfrac{2x-3}{x-1}   \leq 0
%                \item \dfrac{4x+5}{x+2}   \geq 0
%                \item \dfrac{x-3}{2x+6}   \leq 0
%                \item \dfrac{4-5x}{2x-1}  \geq 0
%                \item 3 - \dfrac{x}{x+2}  \leq 0
%                \item \dfrac{3x-2}{5-2x}  \geq 0
%                \item \dfrac{5x}{x-4}     \geq 10
%                \item \dfrac{2x-7}{x-2}   \geq 3
%                \item \dfrac{3x+10}{2x-5} \geq -3
%                \item \dfrac{3x-4}{1-6x}  \leq 2
%                \item \dfrac{6-x}{x-4}    \geq 1
%            \end{mathlist}
%        \end{spacing}
%        \todo{Copiar a resposta do exercício 1 da pg 130 do arquivo 2 do livro do Francisco}
%    \end{exercicio}
%
%    \begin{exercicio}
%        Resolva as inequações e represente a solução na reta real
%        \begin{spacing}{2}
%            \begin{mathlist}[2]
%                \item \dfrac{x^2 -3x+2}{x^2 +x} \geq 0
%                \item \dfrac{x+8}{x^2 +7x+12} \leq 0
%                \item \dfrac{x^2 +5}{x^2 -4} \geq 5
%                \item \dfrac{x^2 +6}{x^2 +1} \leq 2
%                \item 1 + \dfrac{2}{x+1} \leq \dfrac{2}{x}
%                \item \dfrac{4-2x}{x^2 -4} \leq 3
%                \item \dfrac{x+6}{3x^2 +2} \geq 1
%                \item \dfrac{x-5}{2x-5} \geq x
%                \item \dfrac{4x-7}{x+2} \leq x - 2
%                \item \dfrac{3x-1}{x+4} + \dfrac{x}{x-5} \leq 0
%                \item \dfrac{x}{x+1} - \dfrac{1}{x-3} \geq 2
%                \item \dfrac{x^2 +2x+3}{x+15} \leq 1
%                \item \dfrac{3x^2 +2x-13}{x^2 -3x-10} \geq 2
%            \end{mathlist}
%        \end{spacing}
%        \todo{Copiar a resposta do exercício 2 da pg 130 do arquivo 2 do livro do Francisco}
%    \end{exercicio}
%
%    \begin{exercicio}
%        Resolva as inequações e represente a solução na reta real
%            \begin{mathlist}
%                \item \sqrt{x} - 8                     \leq 0
%                \item \sqrt{x} - 3                     \geq 0
%                \item \sqrt{x} + 10                    \geq 0
%                \item \sqrt{6 - 5x} - 4                \geq 0
%                \item \sqrt{2x - 3}                    \leq 5
%                \item \sqrt{3x + 12} + 7               \geq 0
%                \item \sqrt{x + 2}                     \geq x - 4
%                \item \sqrt{10 - 3x}                   \leq x - 2
%                \item \sqrt{8x + 9}                    \leq 2x + 1
%                \item \sqrt{9x^2 - 1}                  \geq 2 - 3x
%                \item \sqrt{8 - 4x^2} + x              \leq 3
%                \item \sqrt{x^2 - 4x - 12}             \leq 3x + 2
%                \item \sqrt{2x^2 + 1} - 2x + 1         \geq 0
%                \item \sqrt{-x^2 + 5x + 14} - 2        \leq x
%            \end{mathlist}
%        \todo{Copiar a resposta do exercício 3 da pg 130 do arquivo 2 do livro do Francisco}
%    \end{exercicio}

    \begin{exercicio}
        Calcule $\abs{3x - 10}$ para $x = 2$ e $x = 5$.
    \end{exercicio}
    
    \begin{exercicio}
        Calcule $\abs{7 - x}$ para $x = -7$, $x = 1$, $x = 7$ e $x = 12$.
    \end{exercicio}
    
    \begin{exercicio}
        Elimine o módulo da expressão $\abs{x}/x^2$.
    \end{exercicio}
    
    \begin{exercicio} % Francisco pg 145, ex 1
        Elimine o módulo das expressões.
        \begin{mathlist}[3]
            \item \abs{8}
            \item \abs{-8}
            \item -\abs{-8}
            \item \abs{3-\pi}
            \item \abs{\pi-3}
            \item \abs{\sqrt{8}-4}
            \item \abs{-\frac{10}{5^2}}
            \item \sqrt{\abs{-4}}
        \end{mathlist}	
        \begin{answers}
            \begin{mathlist}[4]
                \item 8
                \item 8
                \item -8
                \item \pi-3
                \item \pi-3
                \item 4-\sqrt{8}
                \item \sfrac{2}{5}
                \item 2
            \end{mathlist}
        \end{answers}
    \end{exercicio}
    
    \begin{exercicio}
        Calcule os valores das expressões.
        \begin{mathlist}[2]
            \item \abs{-6 + 2}
            \item 4 + \abs{-4}
            \item \frac{\abs{-12 + 4}}{\abs{16 - 12}}
            \item \abs{\frac{0,2 - 1,4}{1,6 - 2,4}}
            \item \sqrt{3\,}\; \abs{-2} + 3\abs{-\sqrt{3\,}}
            \item \abs{\pi - 1} + 2
            \item \abs{2\sqrt{3} - 5} - \abs{\sqrt{3} - 4}
        \end{mathlist}
        \begin{answers}
            \begin{mathlist}[4]
                \item 4
                \item 8
                \item 2
                \item \sfrac{1}{4}
                \item 5\sqrt{3\,}
                \item \pi + 1
                \item 1-\sqrt{3\,}
            \end{mathlist}
        \end{answers}
    \end{exercicio}
    
    \begin{exercicio}
        Suponha que $a$ e $b$ são números reais não-nulos e que $a > b$.
        Determine se as desigualdades são verdadeiras ou falsas.
        \begin{mathlist}[2]
            \item b - a > 0
            \item \frac{a}{b} > 1
            \item a^2 > b^2
            \item -a < -b
        \end{mathlist}	
    \end{exercicio}
    
    \begin{exercicio}
        Determine se as afirmações são verdadeiras para quaisquer números reais $a$ e $b$.
        \begin{mathlist}[2]
            \item \abs{-a} = a
            \item \abs{b^2} = b^2
            \item \abs{a - 4} = \abs{4 - a}
            \item \abs{a + 1} = \abs{a} + 1
            \item \abs{a + b} = \abs{a} + \abs{b}
        \end{mathlist}    
    \end{exercicio}
        
    \begin{exercicio}
        Determine se as afirmações são verdadeiras ou falsas.
	Se verdadeiras, explique o porquê. Se falsas, justifique através de um exemplo.
        \begin{alphalist}
            \item Se $a < b$, então $a - c > b - c$
            \item $\abs{a - b} \leq \abs{b} + \abs{a}$
        \end{alphalist}        
    \end{exercicio}
                
    \begin{exercicio} % Exercício 4 da pg 145 do arquivo 2 do livro do Francisco
        Elimine o módulo das expressões escrevendo-as na forma ``por partes''
        \begin{mathlist}[3]
            \item -\abs{x}
            \item \abs{x}-5
            \item 5-\abs{x}
            \item \abs{x-5}
            \item \abs{5-x}
            \item \abs{5x+1}
            \item \abs{4-3x}
            \item \abs{x^2+7}
            \item \abs{x^2-9}
        \end{mathlist}
        \begin{answers}
            \begin{mathlist}
                \item -\abs{x} = \begin{cases}
                    -x, &\text{se } x \geq 0 \\
                    \hphantom{-}x, &\text{se } x < 0
                \end{cases}
                \item \abs{x}-5 = \begin{cases}
                    \hphantom{-}x-5, &\text{se } x \geq 0 \\
                    -x-5, &\text{se } x < 0
                \end{cases}
                \item 5-\abs{x} = \begin{cases}
                    5-x, &\text{se } x \geq 0 \\
                    5+x, &\text{se } x < 0
                \end{cases}
                \item \abs{x-5} = \begin{cases}
                    x-5, &\text{se } x \geq 5 \\
                    5-x, &\text{se } x < 5
                \end{cases}
                \item \abs{5-x} = \begin{cases}
                    x-5, &\text{se } x \geq 5 \\
                    5-x, &\text{se } x < 5
                   \end{cases}
                \item \abs{5x+1} = \begin{cases}
                	\hphantom{-}5x+1, & \text{se } x \geq \sfrac{-1}{5} \\
                	-5x-1,            & \text{se } x <    \sfrac{-1}{5}
                \end{cases}
                \item \abs{4-3x} = \begin{cases}
                    4-3x, &\text{se } x \leq \sfrac{4}{3} \\
                    3x-4, &\text{se } x > \sfrac{4}{3}
                \end{cases}
                \item \abs{x^2+7} = x^2+7
                \item \abs{x^2-9} = \begin{cases}
                    x^2-9, &\text{se } x \leq -3 \text{ ou } x \geq 3 \\
                    9-x^2, &\text{se } -3 < x 3
                \end{cases}
            \end{mathlist}
        \end{answers}
    \end{exercicio}

    \begin{exercicio} % Exercício 6 da pg 145 do arquivo 2 do livro do Francisco
        Calcule as expressões
        \begin{mathlist}[2]
            \item \abs{5\times(-3)}
            \item -3\abs{-5}
            \item \abs{\frac{-3}{-6}}
            \item \abs{\frac{5-17}{15-6}}
            \item \abs{-2}+6\abs{-5}
            \item \abs{\vphantom{\sum}-2+\abs{-5}}
        \end{mathlist}
        \begin{answers}
            \begin{mathlist}[2]
                \item \abs{5\times(-3)} = 15 
                \item -3\abs{-5} = -15 
                \item \abs{\frac{-3}{-6}} = \sfrac{1}{2}
                \item \abs{\frac{5-17}{15-6}} = \sfrac{3}{4}
                \item \abs{-2}+6\abs{-5} = 32
                \item \abs{\vphantom{\sum}-2+\abs{-5}} = 3
           \end{mathlist}
        \end{answers}        
    \end{exercicio}

    \begin{exercicio} % Exercício 8 da pg 145 do arquivo 2 do livro do Francisco
        Calcule as expressões 
        \begin{mathlist}[2]
            \item \abs{(-4x)\times(-6)}
            \item \abs{ \frac{3x}{-6}}
            \item \abs{-\frac{-3}{6x}}
            \item \abs{-4x}+\abs{8x}
            \item \abs{2x}-\abs{-2x}
            \item \abs{\frac{2x}{3}} - \frac{\abs{x}}{6}
            \item \frac{\abs{2x^2}}{\abs{-4xy}}
            \item \dfrac{\sqrt{x^2}}{\abs{x}}
            \item \sqrt{\abs{-2x^2}}
        \end{mathlist}	
        \begin{answers}
            \begin{mathlist}[2]
                \item \abs{(-4x)\times(-6)}=24\abs{x}
                \item \abs{ \frac{3x}{-6}}=\frac{\abs{x}}{2}
                \item \abs{-\frac{-3}{6x}}=\frac{1}{2\abs{x}}
                \item \abs{-4x}+\abs{8x}=12\abs{x}
                \item \abs{2x}-\abs{-2x}=0
                \item \abs{\frac{2x}{3}} - \frac{\abs{x}}{6}=\frac{\abs{x}}{2}
                \item \frac{\abs{2x^2}}{\abs{-4xy}} =\frac{\abs{x}}{2\abs{y}}
                \item \frac{\sqrt{x^2}}{\abs{x}}=1
                \item \sqrt{\abs{-2x^2}} =\abs{x}\sqrt{2}
            \end{mathlist}
        \end{answers}   
    \end{exercicio}
              
    \begin{exercicio} % exercício 7 da pg 145 do arquivo 2 do livro do Francisco
        Calcule a distância entre os pontos da reta real.
        \begin{alphalist}
            \item \(x_1 = -5\) e \(x_2 = -8\)
            \item \(x_1 = -10\) e \(x_2 = 10\)
            \item \(x_1 = 4,7\) e \(x_2 = 1,2\)
            \item \(x_1 = 2\) e \(x_2 = -12\)
        \end{alphalist}
        \begin{answers}
            \begin{mathlist}[4]
                \item 3
                \item 20
                \item 3,5
                \item 14
            \end{mathlist}
        \end{answers}   
    \end{exercicio}
        
    \begin{exercicio}
        Reescreva as frases abaixo usando equações modulares.
        \begin{alphalist}
            \item A distância entre $x$ e $ 2$ é igual a $3$.
            \item A distância entre $s$ e $-3$ é igual a $4$.
            \item A casa de minha avó e a casa de meu tio estão a \SI{5}{\kilo\meter} de distância.
            \item A farmácia e a padaria estão à mesma distância de minha casa.
        \end{alphalist}
    \end{exercicio}

    \begin{exercicio}
        Determine os pontos da reta real que estão a uma distância de $10$ unidades de $6$.
    \end{exercicio}

    \begin{exercicio}
        Determine os pontos da reta real que estão a uma distância de $2/3$ unidades de $-1$.
    \end{exercicio}

    \begin{exercicio}
        Determine o conjunto dos valores de $x$ que satisfazem a inequação $0 < x^2 - x < 20$
    \end{exercicio}

    \begin{exercicio} % exercício 12 da pg 145 do arquivo 2 do livro do Francisco
        Resolva as equações.
        \begin{mathlist}[3]
            \item \abs{x} = 4
            \item \abs{x} = -4
            \item x = \abs{-4}
            \item x = \abs{4}
            \item \abs{x} = \abs{4}
            \item \abs{x} = \abs{-4}
        \end{mathlist}
        \begin{answers}
            \begin{mathlist}[2]
                \item x=-4 \text{ ou } x=4
                \item \) Não há solução
                \item x=4
                \item x=4
                \item x=-4 \text{ ou } x=4 
                \item x=-4 \text{ ou } x=4 
            \end{mathlist}
        \end{answers}
    \end{exercicio}

    \begin{exercicio} % exercício 13 da pg 146 do arquivo 2 do livro do Francisco
        Resolva as equações.
        \begin{mathlist}[2]
            \item \abs{x-3} = 4
            \item \abs{x-\frac{1}{2}} = 2
            \item \abs{4-x} = \frac{1}{10}
            \item \abs{3x-1} =6
            \item \abs{x-2} =-1
            \item \abs{\frac{x-3}{4}} = 12
            \item \abs{5-4x} =1
            \item \abs{5x-2} =13
            \item \abs{\frac{2-3x}{4}} =5
            \item \abs{\frac{3x}{2}-1} =\frac{5}{2}
            \item \abs{5x-3} =3x+15
            \item \abs{2x-3} =5-4x
            \item \abs{6-x} +5x=7
            \item \abs{\frac{7+4x}{3}} =5+x
            \item \abs{3x+5} =\abs{x-3}
            \item \abs{2x+1} = \abs{2-5x}
            \item \abs{x-1} + \abs{x+2} =5
            \item \abs{2x-5} - \abs{x-2} = 3x+1
        \end{mathlist}
        \begin{answers}
            \begin{mathlist}[2]
                \item x=-4 \text{ ou } x=-2 \text{ ou } x=2 \text{ ou } x=4
                \item x=\sfrac{-3}{2} \text{ ou } x=\sfrac{3}{2}
                \item x=-1 \text{ ou } x=7 \text{ ou } x=3-\sqrt{2} \text{ ou } x=3+\sqrt{2}
                \item x=0 \text{ ou } x=4
                \item x=\sfrac{-2}{3} \text{ ou } x=\sfrac{-1}{2} \text{ ou } x=0 \text{ ou } x=2
                \item x=-6 \text{ ou } x=1 \text{ ou } x=2 \text{ ou } x=3
                \item x=3 \text{ ou } x=-3-\sqrt{2}
                \item x=-2 \text{ ou } x=2
                \item x=-7 \text{ ou } x=2 \text{ ou } x=3
                \item x=-\sqrt{13} \text{ ou } x=\sqrt{13}
            \end{mathlist}
        \end{answers}
    \end{exercicio}

%\todo{
%    \begin{exercicio} % exercício 14 da pg 146 do arquivo 2 do livro do Francisco
%        Resolva as equações.
%        \begin{answers}
%            \begin{mathlist}[2]
%                \item 
%            \end{mathlist}
%        \end{answers}
%    \end{exercicio}
%
%    \begin{exercicio} % exercício 15 da pg 146 do arquivo 2 do livro do Francisco
%        Resolva as equações.
%        \begin{answers}
%            \begin{mathlist}[2]
%                \item 
%            \end{mathlist}
%        \end{answers}
%    \end{exercicio}
%
%    \begin{exercicio} % exercício 17 da pg 146 do arquivo 2 do livro do Francisco
%        Resolva as desigualdades.
%        \begin{answers}
%            \begin{mathlist}[2]
%                \item 
%            \end{mathlist}
%        \end{answers}
%    \end{exercicio}
%
%    \begin{exercicio} % exercício 18 da pg 146 do arquivo 2 do livro do Francisco
%        Identifique, na reta real, os intervalos definidos pelas desigualdades.
%        \begin{answers}
%            \begin{mathlist}[2]
%                \item 
%            \end{mathlist}
%        \end{answers}
%    \end{exercicio}
%
%    \begin{exercicio} % exercício 19 da pg 146 do arquivo 2 do livro do Francisco
%       Resolva as desigualdades.
%       \begin{answers}
%           \begin{mathlist}[2]
%               \item 
%            \end{mathlist}
%        \end{answers}
%    \end{exercicio}
%}
    
\end{exercicios}

%------------------------------------------------------------------------------%
